\documentclass[10pt,a4paper]{article}

\usepackage[T1]{fontenc}
\usepackage[utf8]{inputenc}
\usepackage{lmodern}
\usepackage{amsmath,amssymb,amsthm}
\usepackage{booktabs,array}
\usepackage{geometry}
\usepackage{xcolor}
\usepackage{enumitem}
\usepackage{microtype}
\usepackage{hyperref}
\usepackage{fancyhdr}

\geometry{margin=1in}
\hypersetup{
  colorlinks=true,
  linkcolor=blue!55!black,
  citecolor=blue!55!black,
  urlcolor=blue!55!black,
  pdftitle={A Bit Too Approximate: A Universal Zero-Vector Forgery against Higher-Bit Approximate-ISIS},
  pdfauthor={Martin Feussner},
  pdfsubject={Verified attack report for IACR ePrint 2022/1144},
  pdfkeywords={digital signatures, lattices, Approx-ISIS, universal forgery, cryptanalysis}
}
\setlist{nosep,leftmargin=*}
\setlength{\parindent}{1em}
\setlength{\parskip}{0.25em}
\pagestyle{fancy}
\fancyhf{}
\fancyhead[L]{A Bit Too Approximate}
\fancyhead[R]{14 October 2026}
\fancyfoot[C]{\thepage}
\renewcommand{\headrulewidth}{0.3pt}

\newtheorem{theorem}{Theorem}
\newtheorem{lemma}{Lemma}
\newcommand{\Z}{\mathbb{Z}}
\newcommand{\norm}[1]{\left\lVert #1\right\rVert_2}

\title{\bfseries A Bit Too Approximate:\\
A Universal Zero-Vector Forgery against Higher-Bit Approximate-ISIS}
\author{Martin Feussner\\
\small Selmer Center, University of Bergen\\
\small \texttt{martin.feussner@uib.no}}
\date{14 October 2026}

\begin{document}
\maketitle

\noindent\fcolorbox{black!50}{black!4}{\parbox{0.94\textwidth}{
This audit used the OpenAI model Daybreak Blue at the ultra reasoning level.
OpenAI Codex generated and selected this paper's title and intentionally chose
``A Bit Too Approximate'' as wordplay on the Higher-Bit Approximate-ISIS
scheme name and on the verifier radius becoming too permissive. OpenAI Codex
performed AI-assisted cryptanalysis, implementation, experimentation, and
technical writing. A separately tasked AI agent performed hostile adversarial
review, and a fresh-context AI agent independently reproduced the result from
the pinned paper. The final exact-target attack-history recheck and independent
Publication Reviewer gate passed, and the package is classified as
\texttt{VERIFIED\_ATTACK}. No independent human verification was performed or
required.}}

\begin{abstract}
We analyze Construction~1 of Le Dévéhat, Shizuya, and Hasegawa's higher-bit
Approximate-ISIS signature proposal, IACR ePrint 2022/1144. Its ordinary
verifier accepts $x$ when
\[
  \norm{x}\leq\beta
  \quad\text{and}\quad
  \norm{b^dAx-H(m)}\leq\alpha.
\]
Section~6.4 says parameter selection must respect
$\norm{x}\leq s$ and
$\norm{e}\leq b^\ell\sigma+\sqrt n\,b^d s$. Hence every threshold completion
following that prescription has
$\alpha\geq\sqrt n\,b^d s$. For all three Figure~3 columns labelled
``This work,'' $b^d s>b^k-1$, while the paper defines
$k=\lceil\log_bq\rceil$ and therefore $q\leq b^k$. The required error radius
already exceeds $\sqrt n(q-1)$, an upper bound on the Euclidean norm of every
canonical vector in $\Z_q^n$. At $x=0$ the public-matrix term vanishes, so the
same public zero vector is accepted for every message and every public key,
with no signing query or secret information.

Full-dimension paper-level experiments covered all three advertised rows, 64
fresh messages per row, and both canonical and centered residue lifts. All
384 target checks accepted on one worker in 11.40 seconds with 61,072 KiB peak
resident memory. Precommitted in-domain diagnostics exercised the matrix,
hash, norm, and comparison paths. The paper does not publish concrete
$\alpha$, $\beta$, or $\sigma$ values, a lift, hash serialization, complete
Ring mapping, or its cited private implementation. We therefore claim neither
an implementation replay nor honest signing. The result is quantified only
over Construction~1 threshold completions that obey Section~6.4.
\end{abstract}

\section{Introduction}

Approximate trapdoors permit a hash-and-sign construction to return a short
preimage whose image is close, rather than equal, to a message-dependent
target. Chen, Genise, and Mukherjee introduced the Approx-ISIS formulation and
the corresponding approximate-trapdoor signature framework~\cite{cgm}. Le
Dévéhat, Shizuya, and Hasegawa later proposed retaining only higher bits of the
public matrix to reduce key and signature sizes~\cite{hbaisis}. Their
Construction~1 follows the earlier hash-and-sign pattern while adding a
$b^d$ approximation factor.

Approximation must remain smaller than the target space. Chen, Genise, and
Mukherjee explicitly identify the generic zero-vector boundary: once the
allowed error covers the target, the zero preimage is a trivial solution. We
make no novelty claim for that observation. The issue here is its concrete
interaction with the later paper's own parameter prescription. Section~6.4
adds a term $\sqrt n\,b^d s$ to the required error bound. At every advertised
``This work'' row, the scalar $b^d s$ alone exceeds the largest canonical
coordinate $q-1$ by a factor between 3.89 and 67.63.

The resulting attack is a universal fresh-message forgery:
\begin{itemize}
  \item choose any message that has not been submitted to a signing oracle;
  \item output the length-$m$ all-zero vector; and
  \item make no signing query and use no secret-key information.
\end{itemize}
It applies pointwise to every public matrix and every random-oracle output.
The experiment is useful for confirming the row transcription, full
dimensions, modular arithmetic, both natural lift conventions, and live
verification paths; the security conclusion itself is algebraic.

\paragraph{Review status.}
Primary reproduction, hostile adversarial review, independent fresh-context
reproduction, the final exact-target attack-history recheck, and independent
Publication Review passed. All required publication gates are complete, and
the result is classified as \texttt{VERIFIED\_ATTACK}.

\section{Target and claim boundary}

\subsection{Pinned target}

The target is the PDF of \emph{On the Higher-bit Version of Approximate
Inhomogeneous Short Integer Solution Problem}, IACR ePrint 2022/1144, by
Anaëlle Le Dévéhat, Hiroki Shizuya, and Shingo Hasegawa~\cite{hbaisis}. The
pinned PDF has SHA-256 digest
\begin{center}
\small\ttfamily
cbbbb33cacef911548d52de7990eff74479dbb2e4ab9467974770107458cabeb.
\end{center}
Its text extraction has SHA-256 digest
\begin{center}
\small\ttfamily
b5b60db7b57c6a0cee37cd0a49a2a5039d8ed0cb7d7ee9ba3c9a19ae909a5bef.
\end{center}
The audit uses the mathematical construction and the three columns marked
``This work'' in Figure~3.

\subsection{Ordinary verification relation}

Construction~1 models
$H:\{0,1\}^*\rightarrow\Z_q^n$ as a random oracle. The public key is a matrix
\[
  A\in\Z_{q/b^d}^{n\times m},
\]
and a signature is a vector $x\in\Z_{q/b^d}^{m}$. The verifier accepts exactly
when
\begin{equation}
  \norm{x}\leq\beta,
  \qquad
  \norm{b^dAx-H(m)}\leq\alpha.
  \label{eq:verify}
\end{equation}
The paper uses the Euclidean norm by default. It does not specify which
integer representative of the second modular vector is normed. We therefore
analyze both common fixed conventions: canonical representatives in
$\{0,\ldots,q-1\}$ and centered representatives around zero. Definition~8's
existential error vector is no stricter than taking a centered representative.

Section~6.4 derives the signing bounds
\begin{equation}
  \norm{x}\leq s,
  \qquad
  \norm{e}\leq b^\ell\sigma+\sqrt n\,b^d s,
  \label{eq:bounds}
\end{equation}
and immediately states that these bounds must be respected when setting the
underlying parameters $\alpha$ and $\beta$. We formalize a
\emph{Section~6.4-compliant completion} as any choice satisfying
\begin{equation}
  \beta\geq s,
  \qquad
  \alpha\geq b^\ell\sigma+\sqrt n\,b^d s.
  \label{eq:completion}
\end{equation}
The claim in this paper is limited to such completions. It does not silently
replace $\alpha$ or $\beta$ with empirical norm entries from Figure~3.

\subsection{Unavailable concrete profile}

Figure~3 prints $\tau$, described as the Gaussian width of the trapdoor $R$,
rather than the $\sigma$ appearing in Equation~\eqref{eq:bounds}. It gives no
concrete $\alpha$ or $\beta$, threshold-rounding rule, residue lift,
hash-to-vector mapping, byte order, key encoding, or signature encoding. The
proof-of-concept source is attributed to private communication. The table
also relies on a Ring implementation without defining a complete mapping from
the matrix construction to that structured representation.

These omissions prevent an exact implementation replay. They do not create
an ambiguity in the quantified claim: increasing either radius preserves
acceptance, $b^\ell\sigma$ is nonnegative, and the conservative lower term
$\sqrt n\,b^d s$ already covers the target space.

\section{Universal zero-vector forgery}

\begin{lemma}[Target-space bounds]
For $u\in\Z_q^n$, the Euclidean norm of its canonical representative is at
most $\sqrt n(q-1)$. The norm of its centered representative is at most
$\sqrt n\lfloor q/2\rfloor$.
\end{lemma}

\begin{proof}
Every canonical coordinate lies between $0$ and $q-1$, so the sum of squared
coordinates is at most $n(q-1)^2$. Every centered coordinate has absolute
value at most $\lfloor q/2\rfloor$. Taking square roots gives the bounds.
\end{proof}

\begin{theorem}[Universal forgery]
Suppose Construction~1 uses a Section~6.4-compliant threshold completion and
a parameter row satisfying $b^d s>b^k-1$, where
$k=\lceil\log_bq\rceil$. Then the length-$m$ vector $x=0$ is accepted for
every message and every public matrix under either the canonical or centered
lift.
\end{theorem}

\begin{proof}
The definition of $k$ implies $q\leq b^k$. From
Equation~\eqref{eq:completion} and $\sigma\geq0$,
\[
 \alpha\geq \sqrt n\,b^d s
          > \sqrt n(b^k-1)
          \geq \sqrt n(q-1).
\]
At $x=0$, $\norm{x}=0\leq\beta$ and $b^dAx=0$ for every $A$. The second
verification value is therefore the norm of $-H(m)$ in $\Z_q^n$. The lemma
places its canonical norm at most at $\sqrt n(q-1)$; the centered bound is no
larger. Both inequalities in Equation~\eqref{eq:verify} hold for every
$H(m)$.
\end{proof}

The attack algorithm is constant: output $0^m$. It needs the public parameter
length $m$ but need not inspect $A$, query a signer, or know a secret key. The
message may be selected after the public key is generated and can be any
fresh attacker-chosen string. This is an efficient existential forgery and,
because the same signature works for every message, a universal forgery.

\section{Affected advertised rows}

Table~\ref{tab:rows} transcribes the three Figure~3 columns marked ``This
work.'' Algorithm~2 states that its optimization requires $q=b^k$, giving
$q=65536$ in each executable row. The proof above needs only the weaker
relation $q\leq b^k$.

\begin{table}[ht]
\centering
\small
\begin{tabular}{@{}lrrrrrrrr@{}}
\toprule
Row & $n$ & $b$ & $k$ & $\ell=d$ & $m$ & $s$ & $b^d s$ & $b^d s/(b^k-1)$\\
\midrule
75.0-bit AISIS  & 512  & 2 & 16 & 11 & 3584 & 1453.0 & 2,975,744.0 & 45.4069\\
155.4-bit AISIS & 1024 & 2 & 16 & 11 & 7168 & 2163.9 & 4,431,667.2 & 67.6229\\
153.9-bit AISIS & 1024 & 4 & 8  & 3  & 7168 & 3989.3 & 255,315.2 & 3.8959\\
\bottomrule
\end{tabular}
\caption{Every advertised ``This work'' row satisfies the theorem by a wide
margin. The claimed AISIS bit costs are labels from Figure~3, not security
estimates produced by this audit.}
\label{tab:rows}
\end{table}

The minimum scalar margin is about $3.8959$. Therefore no reasonable rounding
of the one-decimal $s$ entries can close the inequality. The additive
$b^\ell\sigma$ term, omitted from the conservative calculation, can only
increase $\alpha$.

\section{Executable reproduction}

\subsection{Design}

The public reproducer is written solely with the Python standard library. For
each affected row it:
\begin{enumerate}
  \item generates a fresh uniform matrix at the full advertised
        $n\times m$ dimensions in $\Z_{q/b^d}$;
  \item derives multiple distinct fresh messages and maps each to
        $\Z_q^n$ through a domain-separated SHAKE256 random-oracle adapter;
  \item evaluates $b^dAx$ with a general matrix-vector routine on the
        length-$m$ zero signature, with no special zero shortcut;
  \item invokes the two inequalities of the ordinary verifier under both
        canonical and centered lifts; and
  \item records machine-readable row parameters, matrix and message-set
        digests, exact norm comparisons, query counts, and diagnostics.
\end{enumerate}
The product for an unchanged matrix and signature is reused across messages.
This is an optimization of the repeated verifier computation, not an
algebraic shortcut: the initial product traverses the complete matrix through
the same general routine also used for the in-domain unit-vector diagnostic.

The SHAKE256 adapter is not presented as an author encoding. Since the proof
is pointwise for every target vector, no choice of random-oracle serialization
can turn a Section~6.4-compliant threshold into a rejecting one. Likewise,
the experiment's uniform public matrices are not trapdoor-certified keys, but
$A0=0$ for every possible matrix. The reproducer therefore tests the claimed
ordinary verifier relation without claiming key-generation conformance.

\subsection{Independent full-row result}

A fresh-context agent derived its implementation from the pinned paper before
reading the primary experiment. It generated three independent full-size
matrices and 64 messages per row. Table~\ref{tab:results} summarizes the
result.

\begin{table}[ht]
\centering
\small
\begin{tabular}{@{}lrrrl@{}}
\toprule
Row & Messages & Canonical & Centered & Matrix SHA-256 prefix\\
\midrule
TW-512-b2  & 64 & 64/64 & 64/64 & \texttt{f447b22372e1b45e}\ldots\\
TW-1024-b2 & 64 & 64/64 & 64/64 & \texttt{aa0fa985706fef88}\ldots\\
TW-1024-b4 & 64 & 64/64 & 64/64 & \texttt{d9fa77d778166b3e}\ldots\\
\bottomrule
\end{tabular}
\caption{Accepted zero-signature checks in the independent reproduction.}
\label{tab:results}
\end{table}

All 384 target checks accepted. The sequential run used one worker, 11.40
seconds wall time, 8.39 CPU seconds, and 61,072 KiB peak resident memory. It
made zero signing queries and generated no honest signature. The independent
results file has SHA-256 digest
\begin{center}
\small\ttfamily
c350afc1310567ce31e93cd3b4456744fd8cabbe5ea17c34f105062d60533700.
\end{center}

\subsection{Precommitted diagnostics}

The experiment does not use an invalid over-$\beta$ vector or a threshold
chosen from an observed sample. Its matrix-path diagnostic evaluates the
well-formed in-domain vector $e_0$ and confirms that $Ae_0$ differs from
$A0$. Distinct-message checks require every oracle vector and multiple
residual norms to differ.

Figure~3's printed empirical $\norm{e}$ values are preregistered as diagnostic
radii only, never as $\alpha$. In the independent 64-message run, the same
well-formed zero signature produced the rejection counts in
Table~\ref{tab:diagnostics}. These rejections show that the message hashing,
lift, norm, and comparator paths are capable of rejecting in-domain inputs.

\begin{table}[ht]
\centering
\small
\begin{tabular}{@{}lrr@{}}
\toprule
Row & Canonical rejects & Centered rejects\\
\midrule
TW-512-b2  & 64/64 & 32/64\\
TW-1024-b2 & 64/64 & 22/64\\
TW-1024-b4 & 64/64 & 59/64\\
\bottomrule
\end{tabular}
\caption{In-domain diagnostic rejections at the printed Figure~3 empirical
$\norm{e}$ radius. These values are not treated as verifier thresholds.}
\label{tab:diagnostics}
\end{table}

The polished public reproducer adds a deterministic target-space diagnostic.
It uses the valid vector
\[
  u_{\mathrm{diag}}=(q/2,\ldots,q/2)\in\Z_q^n,
\]
which must reject at the same printed diagnostic radius under both lifts. This
vector and radius are fixed before any sample is drawn. It avoids a flaky or
post-selected check while exercising the same residual-norm comparison.

\section{Specification inconsistencies and limitations}

The following observations bound the result and should be resolved by any
future specification:
\begin{enumerate}
  \item \textbf{Missing thresholds.} Figure~3 publishes no concrete
        $\alpha$, $\beta$, or $\sigma$. The result covers every completion
        obeying Section~6.4 but does not describe an unavailable alternative
        profile that ignores that section.
  \item \textbf{Missing representation.} The paper specifies neither a
        modular lift nor hash, key, or signature serialization. Both natural
        fixed lifts pass; the proof holds for every hash vector.
  \item \textbf{Private implementation.} The proof-of-concept is cited as
        private communication. No source replay or interoperability claim is
        made.
  \item \textbf{Honest signing.} The audit did not attempt honest acceptance.
        Doing so would require inventing missing $\alpha$, $\beta$, $\sigma$,
        $\chi$, Ring mapping, and serialization details. The universal attack
        is independent of an honest signing distribution.
  \item \textbf{Row-3 contradiction.} Figure~3 reports $s=3989.3$ but
        $\norm{y}=50533.7$ in the third ``This work'' row, contrary to a
        literal reading of Section~6.4's $\norm{x}\leq s$ bound.
  \item \textbf{Size mismatch.} Figure~2's full matrix public-key size formula
        does not reproduce Figure~3's kilobyte figures without an unstated
        structured representation.
  \item \textbf{Empirical error norms.} The printed $\norm{e}$ values are far
        below the Section~6.4 worst-case bound and are not identified as
        $\alpha$. They cannot be silently substituted for the verifier
        threshold.
\end{enumerate}

The result is a fresh-message forgery, not secret-key or trapdoor recovery.
It says nothing about a future revision that changes the correctness analysis
and threshold rule.

\section{Prior work, attack history, and repair direction}

Chen, Genise, and Mukherjee introduced Approx-ISIS and approximate trapdoors
for smaller hash-and-sign signatures~\cite{cgm}. The zero-vector boundary is
part of that generic Approx-ISIS context, and this report claims no novelty
for it. The concrete finding is that the later higher-bit proposal's explicit
Section~6.4 prescription crosses that boundary at each advertised row.

No prior public attack against the Higher-Bit Approximate-ISIS hash-and-sign
signature of ePrint 2022/1144 was located in the searches performed as of
2026-10-14. This is a dated, bounded result and cannot exclude private,
deleted, or unindexed work; it is not an absolute priority claim. The final
search confidence is moderate-high.

A repair must make the verifier's accepted error region meaningfully smaller
than the target space while retaining an honest-signature correctness margin.
That requires revisiting the higher-bit approximation, signing-error analysis,
and concrete thresholds together. Merely changing encodings, choosing a
different lift, or rejecting the literal all-zero byte string has not been
shown to restore unforgeability: a sound revision needs a complete
specification and a new proof and full-parameter evaluation. This report does
not validate a drop-in patch.

\section{Reproducibility and provenance}

The public artifact contains the complete standard-library source,
three-row parameter transcription, deterministic quick and 64-message
reference outputs, build files, this \LaTeX{} source and PDF, and a SHA-256
manifest. Ordinary invocations draw a fresh 256-bit operating-system seed;
reference invocations use disclosed deterministic fixtures. Every run is
sequential with one worker and emits a machine-readable JSON record ending in
the semantic banner \texttt{ATTACK REPRODUCTION: PASS} on success.

The technical path was checked in three roles: primary reproduction, hostile
adversarial review, and fresh-context independent reproduction. The final
exact-target history recheck and independent Publication Review also passed.
The package retains the stricter claim boundary required by those reviews and
is classified as \texttt{VERIFIED\_ATTACK}. No independent human verification
was performed or required.

The final-history handoff has SHA-256 digest
\begin{center}
\small\ttfamily
daadbc0743cb5ff1c76d2b1cd842261157e3dd438f06bc35579ee1529431c81b,
\end{center}
and its sealed evidence manifest has SHA-256 digest
\begin{center}
\small\ttfamily
9702681ceaaeb90882e0d04c9058f4558fe3f0a5412be2211bd3ede6212d4657.
\end{center}

\section*{Acknowledgements}
The computations were performed on the Norwegian Research and Education Cloud (NREC), using resources provided by the University of Bergen and the University of Oslo.

\begin{thebibliography}{9}

\bibitem{hbaisis}
Anaëlle Le Dévéhat, Hiroki Shizuya, and Shingo Hasegawa.
\newblock On the Higher-bit Version of Approximate Inhomogeneous Short Integer
Solution Problem.
\newblock IACR Cryptology ePrint Archive, Report 2022/1144, 2022.
\newblock \url{https://eprint.iacr.org/2022/1144}.

\bibitem{cgm}
Yilei Chen, Nicholas Genise, and Pratyay Mukherjee.
\newblock Approximate Trapdoors for Lattices and Smaller Hash-and-Sign
Signatures.
\newblock In \emph{Advances in Cryptology--ASIACRYPT 2019}, LNCS 11923,
pages 3--32. Springer, 2019.
\newblock \url{https://doi.org/10.1007/978-3-030-34618-8_1}.

\end{thebibliography}

\end{document}

